% Pollis Domain Packages: SmartRegions.jl
% Compile from this folder with: julia -e 'using PollisTeX; display(compile("SmartRegions.tex"))'
% The values of the first table come from values.jl, read live on the day the text names; those of the second from the panels' Example Code, run the same day.
% The look of a Pollis method vignette (the fonts and colours of the Pollis New Methods documents).
% Set \vignetteauthor, \vignettedate and, if not CC BY-SA 4.0, \vignettelicense after % The look of a Pollis method vignette (the fonts and colours of the Pollis New Methods documents).
% Set \vignetteauthor, \vignettedate and, if not CC BY-SA 4.0, \vignettelicense after % The look of a Pollis method vignette (the fonts and colours of the Pollis New Methods documents).
% Set \vignetteauthor, \vignettedate and, if not CC BY-SA 4.0, \vignettelicense after \input{vignette_style}.
\documentclass[10pt,a4paper]{article}

\usepackage[a4paper,left=18mm,right=18mm,top=15mm,bottom=17mm,footskip=8mm]{geometry}
\usepackage{iftex}
\ifPDFTeX\usepackage[T1]{fontenc}\fi
\usepackage[default,semibold]{sourcesanspro}
\usepackage{textcomp}
\usepackage[italic]{mathastext}
\usepackage{amsmath}
\usepackage{microtype}
\usepackage[table]{xcolor}
\usepackage{tikz}
\usepackage[most]{tcolorbox}
\usepackage{fontawesome5}
\usepackage{tabularx,booktabs}
\usepackage{enumitem}
\usepackage{fancyhdr}
\usepackage[hidelinks]{hyperref}

% ---------------------------------------------------------------- colours (as in the brochure)
\definecolor{navy}{HTML}{1C3A5E}
\definecolor{pblue}{HTML}{3B78B5}
\definecolor{amber}{HTML}{F2B33D}
\definecolor{deepamber}{HTML}{D9902A}
\definecolor{pale}{HTML}{EEF3F9}
\definecolor{ink}{HTML}{2A2F36}
\definecolor{logonavy}{HTML}{173D63}
\definecolor{logodark}{HTML}{C98A1E}
\definecolor{logoamber}{HTML}{F2B446}
\definecolor{logocream}{HTML}{FEEFD2}
\definecolor{soft}{HTML}{5E6773}
\definecolor{agree}{HTML}{2F7D4F}
\definecolor{differ}{HTML}{B23B3B}
\color{ink}

\setlength{\parindent}{0pt}
\setlength{\parskip}{4pt}
\setlist{nosep,leftmargin=*,itemsep=2pt}
\renewcommand{\arraystretch}{1.15}

\pagestyle{fancy}
\fancyhf{}
\renewcommand{\headrulewidth}{0pt}
\fancyfoot[L]{\footnotesize\color{soft}\textbf{Pollis}\enspace New Methods\enspace·\enspace\nmtheme}
\fancyfoot[R]{\footnotesize\color{soft}\thepage}
\newcommand{\nmtheme}{}
\newcommand{\vignetteauthor}{}
\newcommand{\vignettedate}{\today}
\newcommand{\vignettelicense}{\href{https://creativecommons.org/licenses/by-sa/4.0/}{CC BY-SA 4.0}}

% The Pollis logo, redrawn so it prints sharply at any size
\newcommand{\pollislogo}[1]{%
  \begin{tikzpicture}[x=#1/512,y=-#1/512,baseline=(current bounding box.center)]
    \fill[logonavy] (0,0) rectangle (512,512);
    \fill[logodark] (36,422) -- (104,256) -- (172,422) -- cycle;
    \fill[logoamber] (104,422) -- (203,180) -- (302,422) -- cycle;
    \fill[logocream] (203,422) -- (339,89) -- (476,422) -- cycle;
  \end{tikzpicture}}

% Top of the first page: logo, name; then the kicker, title and a one-line promise
\newcommand{\nmhead}[3]{%
  \renewcommand{\nmtheme}{#1}%
  \thispagestyle{fancy}%
  \begin{minipage}[c]{13mm}\pollislogo{11mm}\end{minipage}%
  \begin{minipage}[c]{\dimexpr\linewidth-13mm\relax}
    {\color{navy}\Large\bfseries Pollis}\enspace{\color{soft}New Methods}\hfill
    {\color{deepamber}\small\bfseries\MakeUppercase{#1}}
  \end{minipage}\par\vspace{7mm}
  {\color{navy}\fontsize{22}{26}\selectfont\bfseries #2}\par\vspace{4pt}
  {\color{soft}\large #3}\par\vspace{5mm}}

\newcommand{\nmsection}[1]{\par\vspace{5pt}{\color{navy}\large\bfseries #1}\par\vspace{1pt}}
\newcommand{\nmsub}[1]{\par\vspace{2pt}{\color{deepamber}\bfseries #1}\par}
\newcommand{\gt}{\,{\color{deepamber}\bfseries ›}\,}
\newcommand{\menupath}[1]{{\bfseries\color{navy}#1}}
\newcommand{\fn}[1]{\texttt{\small #1}}
% Citations: \nmcite{key}{Breiman (2001)} links to the entry \nmref{key}{...} in the References
\newcommand{\nmcite}[2]{\hyperlink{ref:#1}{\color{pblue}#2}}
\newcommand{\nmref}[2]{\raisebox{1.1em}[0pt][0pt]{\hypertarget{ref:#1}{}}#2\par}
\newcommand{\nmdoi}[1]{\href{https://doi.org/#1}{\color{pblue}doi:#1}}
% Where the package is: its name and a link (its GitHub repository, once published)
\newcommand{\package}[2]{\href{#2}{\fn{#1.jl}}}

% Key facts at the top: source, missing in Julia, package, result
\newcommand{\factsheet}[4]{%
  \begin{tcolorbox}[enhanced, colback=pale, colframe=pale, boxrule=0pt, sharp corners,
    left=7pt, right=7pt, top=5pt, bottom=5pt, fontupper=\small]
  \begin{tabularx}{\linewidth}{@{}>{\leavevmode\bfseries\color{navy}}p{27mm}>{\leavevmode}X@{}}
  Tested against & #1\\
  Missing in Julia & #2\\
  Package & #3\\
  Result & #4
  \end{tabularx}
  \end{tcolorbox}}

% Boxes
\newtcolorbox{pbox}[2][]{enhanced, colback=pale, colframe=amber, boxrule=0pt,
  leftrule=3pt, sharp corners, left=7pt, right=7pt, top=6pt, bottom=6pt,
  fontupper=\small\raggedright, before upper={{\color{navy}\normalsize\bfseries #2}\par\smallskip}, #1}
\newtcolorbox{nbox}[2][]{enhanced, colback=navy, colframe=navy, boxrule=0pt,
  sharp corners, left=8pt, right=8pt, top=7pt, bottom=7pt, coltext=white,
  fontupper=\small\raggedright, before upper={{\color{amber}\normalsize\bfseries #2}\par\smallskip}, #1}

% Value tables: published beside Julia
\newenvironment{valuetable}[1]{%
  \par\small\tabularx{\linewidth}{#1}\toprule}{\bottomrule\endtabularx\par}
\newcommand{\ok}{{\color{agree}\faCheck}}
\newcommand{\no}{{\color{differ}\faTimes}}

% End of the document: credit and licence
\newcommand{\nmcredit}{%
  \par\vfill
  \begin{minipage}[b]{0.78\linewidth}\footnotesize\color{soft}
  A Pollis method, tested against the published tables on the published data.\par
  \textcopyright\ \vignetteauthor\enspace·\enspace Licensed under \vignettelicense\enspace·\enspace\vignettedate
  \end{minipage}\hfill
  \begin{minipage}[b]{0.18\linewidth}\raggedleft\pollislogo{9mm}\end{minipage}}
.
\documentclass[10pt,a4paper]{article}

\usepackage[a4paper,left=18mm,right=18mm,top=15mm,bottom=17mm,footskip=8mm]{geometry}
\usepackage{iftex}
\ifPDFTeX\usepackage[T1]{fontenc}\fi
\usepackage[default,semibold]{sourcesanspro}
\usepackage{textcomp}
\usepackage[italic]{mathastext}
\usepackage{amsmath}
\usepackage{microtype}
\usepackage[table]{xcolor}
\usepackage{tikz}
\usepackage[most]{tcolorbox}
\usepackage{fontawesome5}
\usepackage{tabularx,booktabs}
\usepackage{enumitem}
\usepackage{fancyhdr}
\usepackage[hidelinks]{hyperref}

% ---------------------------------------------------------------- colours (as in the brochure)
\definecolor{navy}{HTML}{1C3A5E}
\definecolor{pblue}{HTML}{3B78B5}
\definecolor{amber}{HTML}{F2B33D}
\definecolor{deepamber}{HTML}{D9902A}
\definecolor{pale}{HTML}{EEF3F9}
\definecolor{ink}{HTML}{2A2F36}
\definecolor{logonavy}{HTML}{173D63}
\definecolor{logodark}{HTML}{C98A1E}
\definecolor{logoamber}{HTML}{F2B446}
\definecolor{logocream}{HTML}{FEEFD2}
\definecolor{soft}{HTML}{5E6773}
\definecolor{agree}{HTML}{2F7D4F}
\definecolor{differ}{HTML}{B23B3B}
\color{ink}

\setlength{\parindent}{0pt}
\setlength{\parskip}{4pt}
\setlist{nosep,leftmargin=*,itemsep=2pt}
\renewcommand{\arraystretch}{1.15}

\pagestyle{fancy}
\fancyhf{}
\renewcommand{\headrulewidth}{0pt}
\fancyfoot[L]{\footnotesize\color{soft}\textbf{Pollis}\enspace New Methods\enspace·\enspace\nmtheme}
\fancyfoot[R]{\footnotesize\color{soft}\thepage}
\newcommand{\nmtheme}{}
\newcommand{\vignetteauthor}{}
\newcommand{\vignettedate}{\today}
\newcommand{\vignettelicense}{\href{https://creativecommons.org/licenses/by-sa/4.0/}{CC BY-SA 4.0}}

% The Pollis logo, redrawn so it prints sharply at any size
\newcommand{\pollislogo}[1]{%
  \begin{tikzpicture}[x=#1/512,y=-#1/512,baseline=(current bounding box.center)]
    \fill[logonavy] (0,0) rectangle (512,512);
    \fill[logodark] (36,422) -- (104,256) -- (172,422) -- cycle;
    \fill[logoamber] (104,422) -- (203,180) -- (302,422) -- cycle;
    \fill[logocream] (203,422) -- (339,89) -- (476,422) -- cycle;
  \end{tikzpicture}}

% Top of the first page: logo, name; then the kicker, title and a one-line promise
\newcommand{\nmhead}[3]{%
  \renewcommand{\nmtheme}{#1}%
  \thispagestyle{fancy}%
  \begin{minipage}[c]{13mm}\pollislogo{11mm}\end{minipage}%
  \begin{minipage}[c]{\dimexpr\linewidth-13mm\relax}
    {\color{navy}\Large\bfseries Pollis}\enspace{\color{soft}New Methods}\hfill
    {\color{deepamber}\small\bfseries\MakeUppercase{#1}}
  \end{minipage}\par\vspace{7mm}
  {\color{navy}\fontsize{22}{26}\selectfont\bfseries #2}\par\vspace{4pt}
  {\color{soft}\large #3}\par\vspace{5mm}}

\newcommand{\nmsection}[1]{\par\vspace{5pt}{\color{navy}\large\bfseries #1}\par\vspace{1pt}}
\newcommand{\nmsub}[1]{\par\vspace{2pt}{\color{deepamber}\bfseries #1}\par}
\newcommand{\gt}{\,{\color{deepamber}\bfseries ›}\,}
\newcommand{\menupath}[1]{{\bfseries\color{navy}#1}}
\newcommand{\fn}[1]{\texttt{\small #1}}
% Citations: \nmcite{key}{Breiman (2001)} links to the entry \nmref{key}{...} in the References
\newcommand{\nmcite}[2]{\hyperlink{ref:#1}{\color{pblue}#2}}
\newcommand{\nmref}[2]{\raisebox{1.1em}[0pt][0pt]{\hypertarget{ref:#1}{}}#2\par}
\newcommand{\nmdoi}[1]{\href{https://doi.org/#1}{\color{pblue}doi:#1}}
% Where the package is: its name and a link (its GitHub repository, once published)
\newcommand{\package}[2]{\href{#2}{\fn{#1.jl}}}

% Key facts at the top: source, missing in Julia, package, result
\newcommand{\factsheet}[4]{%
  \begin{tcolorbox}[enhanced, colback=pale, colframe=pale, boxrule=0pt, sharp corners,
    left=7pt, right=7pt, top=5pt, bottom=5pt, fontupper=\small]
  \begin{tabularx}{\linewidth}{@{}>{\leavevmode\bfseries\color{navy}}p{27mm}>{\leavevmode}X@{}}
  Tested against & #1\\
  Missing in Julia & #2\\
  Package & #3\\
  Result & #4
  \end{tabularx}
  \end{tcolorbox}}

% Boxes
\newtcolorbox{pbox}[2][]{enhanced, colback=pale, colframe=amber, boxrule=0pt,
  leftrule=3pt, sharp corners, left=7pt, right=7pt, top=6pt, bottom=6pt,
  fontupper=\small\raggedright, before upper={{\color{navy}\normalsize\bfseries #2}\par\smallskip}, #1}
\newtcolorbox{nbox}[2][]{enhanced, colback=navy, colframe=navy, boxrule=0pt,
  sharp corners, left=8pt, right=8pt, top=7pt, bottom=7pt, coltext=white,
  fontupper=\small\raggedright, before upper={{\color{amber}\normalsize\bfseries #2}\par\smallskip}, #1}

% Value tables: published beside Julia
\newenvironment{valuetable}[1]{%
  \par\small\tabularx{\linewidth}{#1}\toprule}{\bottomrule\endtabularx\par}
\newcommand{\ok}{{\color{agree}\faCheck}}
\newcommand{\no}{{\color{differ}\faTimes}}

% End of the document: credit and licence
\newcommand{\nmcredit}{%
  \par\vfill
  \begin{minipage}[b]{0.78\linewidth}\footnotesize\color{soft}
  A Pollis method, tested against the published tables on the published data.\par
  \textcopyright\ \vignetteauthor\enspace·\enspace Licensed under \vignettelicense\enspace·\enspace\vignettedate
  \end{minipage}\hfill
  \begin{minipage}[b]{0.18\linewidth}\raggedleft\pollislogo{9mm}\end{minipage}}
.
\documentclass[10pt,a4paper]{article}

\usepackage[a4paper,left=18mm,right=18mm,top=15mm,bottom=17mm,footskip=8mm]{geometry}
\usepackage{iftex}
\ifPDFTeX\usepackage[T1]{fontenc}\fi
\usepackage[default,semibold]{sourcesanspro}
\usepackage{textcomp}
\usepackage[italic]{mathastext}
\usepackage{amsmath}
\usepackage{microtype}
\usepackage[table]{xcolor}
\usepackage{tikz}
\usepackage[most]{tcolorbox}
\usepackage{fontawesome5}
\usepackage{tabularx,booktabs}
\usepackage{enumitem}
\usepackage{fancyhdr}
\usepackage[hidelinks]{hyperref}

% ---------------------------------------------------------------- colours (as in the brochure)
\definecolor{navy}{HTML}{1C3A5E}
\definecolor{pblue}{HTML}{3B78B5}
\definecolor{amber}{HTML}{F2B33D}
\definecolor{deepamber}{HTML}{D9902A}
\definecolor{pale}{HTML}{EEF3F9}
\definecolor{ink}{HTML}{2A2F36}
\definecolor{logonavy}{HTML}{173D63}
\definecolor{logodark}{HTML}{C98A1E}
\definecolor{logoamber}{HTML}{F2B446}
\definecolor{logocream}{HTML}{FEEFD2}
\definecolor{soft}{HTML}{5E6773}
\definecolor{agree}{HTML}{2F7D4F}
\definecolor{differ}{HTML}{B23B3B}
\color{ink}

\setlength{\parindent}{0pt}
\setlength{\parskip}{4pt}
\setlist{nosep,leftmargin=*,itemsep=2pt}
\renewcommand{\arraystretch}{1.15}

\pagestyle{fancy}
\fancyhf{}
\renewcommand{\headrulewidth}{0pt}
\fancyfoot[L]{\footnotesize\color{soft}\textbf{Pollis}\enspace New Methods\enspace·\enspace\nmtheme}
\fancyfoot[R]{\footnotesize\color{soft}\thepage}
\newcommand{\nmtheme}{}
\newcommand{\vignetteauthor}{}
\newcommand{\vignettedate}{\today}
\newcommand{\vignettelicense}{\href{https://creativecommons.org/licenses/by-sa/4.0/}{CC BY-SA 4.0}}

% The Pollis logo, redrawn so it prints sharply at any size
\newcommand{\pollislogo}[1]{%
  \begin{tikzpicture}[x=#1/512,y=-#1/512,baseline=(current bounding box.center)]
    \fill[logonavy] (0,0) rectangle (512,512);
    \fill[logodark] (36,422) -- (104,256) -- (172,422) -- cycle;
    \fill[logoamber] (104,422) -- (203,180) -- (302,422) -- cycle;
    \fill[logocream] (203,422) -- (339,89) -- (476,422) -- cycle;
  \end{tikzpicture}}

% Top of the first page: logo, name; then the kicker, title and a one-line promise
\newcommand{\nmhead}[3]{%
  \renewcommand{\nmtheme}{#1}%
  \thispagestyle{fancy}%
  \begin{minipage}[c]{13mm}\pollislogo{11mm}\end{minipage}%
  \begin{minipage}[c]{\dimexpr\linewidth-13mm\relax}
    {\color{navy}\Large\bfseries Pollis}\enspace{\color{soft}New Methods}\hfill
    {\color{deepamber}\small\bfseries\MakeUppercase{#1}}
  \end{minipage}\par\vspace{7mm}
  {\color{navy}\fontsize{22}{26}\selectfont\bfseries #2}\par\vspace{4pt}
  {\color{soft}\large #3}\par\vspace{5mm}}

\newcommand{\nmsection}[1]{\par\vspace{5pt}{\color{navy}\large\bfseries #1}\par\vspace{1pt}}
\newcommand{\nmsub}[1]{\par\vspace{2pt}{\color{deepamber}\bfseries #1}\par}
\newcommand{\gt}{\,{\color{deepamber}\bfseries ›}\,}
\newcommand{\menupath}[1]{{\bfseries\color{navy}#1}}
\newcommand{\fn}[1]{\texttt{\small #1}}
% Citations: \nmcite{key}{Breiman (2001)} links to the entry \nmref{key}{...} in the References
\newcommand{\nmcite}[2]{\hyperlink{ref:#1}{\color{pblue}#2}}
\newcommand{\nmref}[2]{\raisebox{1.1em}[0pt][0pt]{\hypertarget{ref:#1}{}}#2\par}
\newcommand{\nmdoi}[1]{\href{https://doi.org/#1}{\color{pblue}doi:#1}}
% Where the package is: its name and a link (its GitHub repository, once published)
\newcommand{\package}[2]{\href{#2}{\fn{#1.jl}}}

% Key facts at the top: source, missing in Julia, package, result
\newcommand{\factsheet}[4]{%
  \begin{tcolorbox}[enhanced, colback=pale, colframe=pale, boxrule=0pt, sharp corners,
    left=7pt, right=7pt, top=5pt, bottom=5pt, fontupper=\small]
  \begin{tabularx}{\linewidth}{@{}>{\leavevmode\bfseries\color{navy}}p{27mm}>{\leavevmode}X@{}}
  Tested against & #1\\
  Missing in Julia & #2\\
  Package & #3\\
  Result & #4
  \end{tabularx}
  \end{tcolorbox}}

% Boxes
\newtcolorbox{pbox}[2][]{enhanced, colback=pale, colframe=amber, boxrule=0pt,
  leftrule=3pt, sharp corners, left=7pt, right=7pt, top=6pt, bottom=6pt,
  fontupper=\small\raggedright, before upper={{\color{navy}\normalsize\bfseries #2}\par\smallskip}, #1}
\newtcolorbox{nbox}[2][]{enhanced, colback=navy, colframe=navy, boxrule=0pt,
  sharp corners, left=8pt, right=8pt, top=7pt, bottom=7pt, coltext=white,
  fontupper=\small\raggedright, before upper={{\color{amber}\normalsize\bfseries #2}\par\smallskip}, #1}

% Value tables: published beside Julia
\newenvironment{valuetable}[1]{%
  \par\small\tabularx{\linewidth}{#1}\toprule}{\bottomrule\endtabularx\par}
\newcommand{\ok}{{\color{agree}\faCheck}}
\newcommand{\no}{{\color{differ}\faTimes}}

% End of the document: credit and licence
\newcommand{\nmcredit}{%
  \par\vfill
  \begin{minipage}[b]{0.78\linewidth}\footnotesize\color{soft}
  A Pollis method, tested against the published tables on the published data.\par
  \textcopyright\ \vignetteauthor\enspace·\enspace Licensed under \vignettelicense\enspace·\enspace\vignettedate
  \end{minipage}\hfill
  \begin{minipage}[b]{0.18\linewidth}\raggedleft\pollislogo{9mm}\end{minipage}}


% A package Pollis needs for panels of its own, not one that reproduces a paper: its own header, key facts and credit
\renewcommand{\vignetteauthor}{Antonio Saragga Seabra}
\renewcommand{\vignettedate}{22 November 2026}
\renewcommand{\vignettelicense}{MIT}
\fancyfoot[L]{\footnotesize\color{soft}\textbf{Pollis}\enspace Domain Packages\enspace·\enspace\nmtheme}
\renewcommand{\nmhead}[3]{%
  \renewcommand{\nmtheme}{#1}%
  \thispagestyle{fancy}%
  \begin{minipage}[c]{13mm}\pollislogo{11mm}\end{minipage}%
  \begin{minipage}[c]{\dimexpr\linewidth-13mm\relax}
    {\color{navy}\Large\bfseries Pollis}\enspace{\color{soft}Domain Packages}\hfill
    {\color{deepamber}\small\bfseries\MakeUppercase{#1}}
  \end{minipage}\par\vspace{7mm}
  {\color{navy}\fontsize{22}{26}\selectfont\bfseries #2}\par\vspace{4pt}
  {\color{soft}\large #3}\par\vspace{5mm}}
\renewcommand{\factsheet}[4]{%
  \begin{tcolorbox}[enhanced, colback=pale, colframe=pale, boxrule=0pt, sharp corners,
    left=7pt, right=7pt, top=5pt, bottom=5pt, fontupper=\small]
  \begin{tabularx}{\linewidth}{@{}>{\leavevmode\bfseries\color{navy}}p{27mm}>{\leavevmode}X@{}}
  In Pollis & #1\\
  Missing in Julia & #2\\
  Package & #3\\
  Tests & #4
  \end{tabularx}
  \end{tcolorbox}}
\renewcommand{\nmcredit}{%
  \par\vfill
  \begin{minipage}[b]{0.78\linewidth}\footnotesize\color{soft}
  Package written and tested with Claude Code (Anthropic) in Pollis.\par
  \textcopyright\ 2026 \vignetteauthor\enspace·\enspace Licensed under \vignettelicense\enspace·\enspace\vignettedate
  \end{minipage}\hfill
  \begin{minipage}[b]{0.18\linewidth}\raggedleft\pollislogo{9mm}\par\vspace{2pt}{\footnotesize contact@pollis.org}\end{minipage}}
\newcolumntype{R}{>{\raggedleft\arraybackslash}X}
\newcolumntype{L}{>{\raggedright\arraybackslash}X}
\newcolumntype{P}[1]{>{\raggedright\arraybackslash}p{#1}}

\begin{document}
\emergencystretch=3em

\nmhead{Smart Regions}
  {SmartRegions: open data of places in Julia}
  {Overture Maps, the OpenSky Network, Mapillary, Copernicus, USGS EarthExplorer, OpenTopography, OpenStreetMap and
   its street networks, Eurostat, INSPIRE and the flows of mobility, read into DuckDB}

\factsheet
  {Powers 14 panels of the Smart Regions Toolbox: OpenStreetMap Data, Overture Maps Data, Eurostat GISCO, INSPIRE Data,
   Copernicus Data Space, USGS EarthExplorer, OpenTopography, Mapillary Street-Level Imagery and Viewer, Route and Network
   Analysis, Mobility and Flows, Air Traffic, Google Maps and Google Photorealistic 3D Maps.}
  {Of the packages of the General registry, by name and description, none reads Overture Maps, the OpenSky Network,
   Mapillary, the Machine-to-Machine API of USGS, OpenTopography, GISCO, Eurostat's API, the services of INSPIRE or
   Geofabrik. For OpenStreetMap the registry has LightOSM, OpenStreetMapX and OpenStreetMapPBF, and for catalogues of
   scenes STAC and CopernicusData; none of them was run here.}
  {\package{SmartRegions}{https://github.com/Trumpingtons/SmartRegions.jl} 0.3.0, MIT:\newline
   \fn{Pkg.Registry.add(url = "https://github.com/Trumpingtons/PollisRegistry.git")}\newline \fn{Pkg.add("SmartRegions")}}
  {\fn{test/runtests.jl}: 13 tests that need no network, among them fastest routes and a matrix of travel times on a
   network of four nodes, against values worked by hand. Every function that reads a source, and every tab of the 14
   panels, run against the live sources on 22 November 2026 (next page).}

\nmsection{Why a package}
Each source of open data has its own way in: a STAC catalogue, a REST API with or without an account, GeoParquet files
over HTTPS, Overpass queries, Atom feeds, Web Feature Services that answer in GML, census tables in CSV, JSON-stat or
Excel. Reading one takes requests, retries, pages, the parsing of XML or JSON, and the conversion of coordinates. The
package does that once, so that a panel's Example Code is a few calls, and leaves what it reads as tables of DuckDB,
which the Next Steps query in SQL and draw on maps. It does not depend on DuckDB.jl: every function takes a connection,
\fn{con}, of the DuckDB the user loads, with the extensions \fn{spatial} and, for files over HTTPS, \fn{httpfs}.

\nmsection{What it reads}
\begin{valuetable}{@{}P{28mm} L P{22mm}@{}}
Source & Functions & Account \\ \midrule
Overture Maps & \fn{overture}, \fn{overture\_release}: the features of a theme over a box or within a radius, only the row groups that overlap it & none \\
OpenSky Network & \fn{opensky\_states}, \fn{opensky\_track}, \fn{opensky\_flights}, \fn{opensky\_credits} & optional client \\
Mapillary & \fn{mapillary\_coverage}, \fn{mapillary\_images}, \fn{mapillary\_near}, \fn{mapillary\_features}, \fn{mapillary\_thumbnails}, \fn{mapillary\_image} & token \\
Copernicus Data Space & \fn{copernicus\_scenes} (STAC), \fn{copernicus\_band} & for bands \\
USGS EarthExplorer & \fn{usgs\_scenes}, \fn{usgs\_download}, \fn{usgs\_bands}, \fn{usgs\_dem} & token \\
OpenTopography & \fn{opentopography\_catalog}, \fn{opentopography\_dem} & key for elevation \\
OpenStreetMap & \fn{osm\_overpass} (now or as of a date), \fn{osm\_read}; \fn{geofabrik\_layer}, \fn{geofabrik\_extract}, \fn{geofabrik\_earlier} & none \\
Street networks & \fn{osm\_streets}, \fn{street\_network}, \fn{nearest\_node}, \fn{fastest\_route}, \fn{travel\_times}, \fn{service\_areas}, \fn{travel\_matrix} & none \\
Eurostat & \fn{gisco\_files}, \fn{gisco\_file} (NUTS, LAU, countries, Urban Audit), \fn{eurostat} (SDMX-CSV), \fn{vector\_units} & none \\
INSPIRE & \fn{inspire\_search}, \fn{inspire\_countries} (data.europa.eu), \fn{wfs\_features}, \fn{atom\_features}, \fn{ogc\_features} & none \\
Mobility & \fn{commuting\_flows} (9 cities), \fn{trip\_flows} (London, New York, Singapore), \fn{traffic\_counts} (5 cities), \fn{air\_flights}, \fn{given\_flows}, \fn{given\_counts} & Singapore: LTA key \\
Google's maps & \fn{googlemap}, \fn{googlemap3d}: a page with one's own features on Google's 2D or 3D map & API key \\
Places & \fn{around}, \fn{utmzone}, \fn{projected\_box}, \fn{geojson} & \\
\end{valuetable}

\noindent Among the sources: OpenStreetMap (\nmcite{haklay2008}{Haklay and Weber, 2008}), the OpenSky Network
(\nmcite{schafer2014}{Schäfer, Strohmeier, Lenders, Martinovic and Wilhelm, 2014}) and its monthly list of flights
(\nmcite{strohmeier2021}{Strohmeier, Olive, Lübbe, Schäfer and Lenders, 2021}), and the services the member states of
the European Union publish under the INSPIRE Directive (\nmcite{inspire2007}{Directive 2007/2/EC}).

\nmsub{Street networks}
\fn{osm\_streets} cuts the ways of OpenStreetMap with a highway tag into segments between nodes, measured on the WGS 84
ellipsoid by DuckDB's \fn{ST\_Distance\_Spheroid}. \fn{street\_network} keeps those a mode may use, by the tags of access:
walking both ways at 4.8 km/h, cycling at 15 km/h with one-way streets one way unless open to bicycles both ways, driving
at the speed limit, else a speed for the class of road. The graph is a \fn{SimpleDiGraph} of Graphs.jl, and routes and
travel times are shortest paths in seconds (\nmcite{dijkstra1959}{Dijkstra, 1959}) from a node of its largest strongly
connected part, in which every node reaches every other.

\nmsub{Files kept}
What is downloaded is kept in a folder, \fn{folder}, and read from there the next time: the files of GISCO, Geofabrik,
the sources of mobility, the answers of Overpass, and the features of WFS, Atom and OGC API services as GeoParquet.

\nmsection{A worked example}
\begin{footnotesize}
\begin{verbatim}
using SmartRegions, DuckDB, Tables
con = DBInterface.connect(DuckDB.DB)
foreach(sql -> DBInterface.execute(con, sql), ["LOAD spatial", "LOAD httpfs", "SET geometry_always_xy = true"])
folder = mkpath("smartregions")

# A walk across central Cambridge, from the station to King's College Chapel
osm_streets(con, (0.105, 52.188, 0.145, 52.216); folder)
net = street_network(con, "walk")
from, _ = nearest_node(con, net, 0.13772, 52.19411)
to, _ = nearest_node(con, net, 0.11654, 52.20478)
path, minutes, km = fastest_route(net, from, to)                  # 28.5 minutes, 2.28 km
travel_times(con, net, from, 15)                                  # 3,054 nodes within 15 minutes

# The NUTS 3 regions of Portugal, with their density of population from Eurostat
nuts = gisco_file("nuts", "NUTS_RG_20M_2024_4326_LEVL_3"; folder)
DBInterface.execute(con, "CREATE TABLE regions AS SELECT * FROM $(vector_units(con, nuts)) WHERE CNTR_CODE = 'PT'")
eurostat(con, "demo_r_d3dens", "A.PER_KM2.")                      # the table eurostat; returns "geo"

# The commuters between the zones of Cambridge, from the Census 2021
commuting_flows(con, "cambridge"; folder)                         # the tables zones and flows
\end{verbatim}
\end{footnotesize}

\nmsection{What the checks show}
\nmsub{Read live with values.jl}
Read on 22 November 2026 by \fn{vignette/values.jl}, from the sources that need no account. A source's answer changes
with the day: the aircraft over Sydney now, the newest release of Overture Maps, the edits of OpenStreetMap.

\begin{valuetable}{@{}P{30mm} L P{34mm}@{}}
Source & Read & Value \\ \midrule
Overture Maps & buildings within 300 m of Martin Place, Sydney; the release & 163; 2026-11-18.0 \\
OpenSky Network & aircraft in a box of 60 km around Sydney Harbour, 01:48 UTC & 33 \\
Copernicus Data Space & Sentinel-2 L2A scenes over central Bologna, June 2025, at most 10\% cloud & 2 \\
OpenTopography & topographic datasets over the same box & 17 \\
Overpass & features tagged \fn{amenity=school}, central Bologna & 54 \\
Geofabrik & points of interest of Monaco & 1,557 \\
Street network & nodes and directed edges open to walking, central Cambridge & 12,570; 27,760 \\
 & fastest walk, the station to King's College Chapel & 28.5 min, 2.28 km \\
 & nodes reached within 15 minutes of the station & 3,054 \\
GISCO and Eurostat & NUTS 3 regions of Portugal; those with a density for 2024 & 26; 26 \\
INSPIRE & datasets of administrative units of the Netherlands in data.europa.eu & 35 \\
 & administrative units by WFS, OGC API and Atom & 353; 353; 355 \\
Commuting & Cambridge, Census 2021: zones; workers between them; to or from outside; at or from home & 33; 56,934; 13,619; 81,824 \\
Trips & London, Santander Cycles, May 2026: zones; journeys between them; to or from outside & 1,002; 833,327; 29,690 \\
Road counts & Cambridge: count points; first and last year & 129; 2000, 2025 \\
Flights & December 2022: flights seen; with both airports & 2,669,349; 1,690,944 \\
\end{valuetable}

\medskip
\newpage
\nmsub{The other tabs of the panels, run the same day}
\begin{valuetable}{@{}P{34mm} P{22mm} L@{}}
Panel & Tab & Result \\ \midrule
OpenStreetMap Data & Overpass & 40 schools in central Bologna as of 1 January 2020 \\
 & Extract & 1,438 amenities of Monaco \\
Route and Network Analysis & Route & a drive in Monaco, from a PBF file of Geofabrik \\
INSPIRE Data & WFS & 340 municipalities of 41,412 km\textsuperscript{2}, beside 342 local administrative units of GISCO, 37,402 km\textsuperscript{2} \\
Mobility and Flows & Commuting & zones: Oxford 87, Amsterdam 49, Paris 142, Frankfurt 27, Tokyo 53, Hong Kong 21, New York 262 \\
 & Trips & New York, September 2026: 3,850,083 taxi trips \\
 & Interaction & Cambridge: distance elasticity $-$1.0818 (SE .0297), common part of commuters .888 \\
 & Counts & count points: New York 3,937, Paris 2,083 \\
 & Air & December 2022: 30,275 routes \\
\end{valuetable}

\medskip
\noindent Singapore's bus trips stop, as they should without a key, with a message that asks for one. The Months Compared
tab was not run: it downloads about 600 MB.

\nmsub{Tests without the network}
On a network of four nodes, 1 to 2 in 60 s over 100 m, 2 to 3 in 120 s over 200 m and 1 to 3 directly in 300 s over
250 m, \fn{fastest\_route} from 1 to 3 goes through 2, in 3 minutes over .3 km, and \fn{travel\_matrix} gives Inf where
no path leads. The UTM zones of Sydney, New York and London are EPSG:32756, 32618 and 32630. A page of Google's maps whose
features hold \fn{</script>} keeps its three scripts whole. Requests the services would refuse stop before they are
sent: a Mapillary search over more than 500 m or 1 square degree, an elevation model of OpenTopography over more than
250 km\textsuperscript{2}, and a city with no source of commuters, trips or counts.

\nmsection{References}
{\footnotesize\raggedright\setlength{\parskip}{1pt}
\nmref{dijkstra1959}{Dijkstra, E. W. (1959). A note on two problems in connexion with graphs. \emph{Numerische Mathematik}, 1, 269--271. \nmdoi{10.1007/BF01386390}}
\nmref{haklay2008}{Haklay, M. and Weber, P. (2008). OpenStreetMap: user-generated street maps. \emph{IEEE Pervasive Computing}, 7(4), 12--18. \nmdoi{10.1109/MPRV.2008.80}}
\nmref{schafer2014}{Schäfer, M., Strohmeier, M., Lenders, V., Martinovic, I. and Wilhelm, M. (2014). Bringing up OpenSky: a large-scale ADS-B sensor network for research. In \emph{Proceedings of the 13th International Symposium on Information Processing in Sensor Networks}, 83--94. \nmdoi{10.1109/IPSN.2014.6846743}}
\nmref{strohmeier2021}{Strohmeier, M., Olive, X., Lübbe, J., Schäfer, M. and Lenders, V. (2021). Crowdsourced air traffic data from the OpenSky Network 2019--2020. \emph{Earth System Science Data}, 13(2), 357--366. \nmdoi{10.5194/essd-13-357-2021}}
\nmref{inspire2007}{Directive 2007/2/EC of the European Parliament and of the Council of 14 March 2007 establishing an Infrastructure for Spatial Information in the European Community (INSPIRE). \emph{Official Journal of the European Union}, L 108, 1--14.}
}

\nmcredit
\end{document}
